% Rutas y nombres en \texttt no se dividen: se deja holgura al espaciado.
\sloppy

\chapter{Mapa de registros}
\label{app:regs}

Este apéndice es el contrato de interfaz entre el motor y su maestro, en
la versión 0x03 del \ac{RTL}. El mismo texto, con la historia de cada
regla, está en el repositorio como \nolinkurl{docs/mapa_registros.md}; lo
siguen el \ac{RTL} (\nolinkurl{rtl/top/motor_top.vhd}), el modelo del PC
(\nolinkurl{firmware/host_check/motor_sim.c}) y la capa de acceso del
firmware (\nolinkurl{firmware/componentes/motor_i2c}).

El motor responde en la dirección \ac{I2C} de 7 bits 0x30. Una escritura
fija un puntero de registro de 8 bits que avanza con cada byte de datos,
leído o escrito. El esclavo nunca estira el reloj: las órdenes largas se
lanzan escribiendo \texttt{CONTROL} y el maestro sondea \texttt{STATUS}
hasta que baja \texttt{ocupado}. Las direcciones no definidas se leen como
0x00 y las escrituras en registros de solo lectura se ignoran.

\begin{table}[htbp]
\centering\small
\caption{Registros. R: lectura, W: escritura.}
\label{tab:app_regs}
\begin{tabular}{@{}llc>{\raggedright\arraybackslash}p{0.54\linewidth}@{}}
\toprule
Dirección & Nombre & Acceso & Contenido \\
\midrule
0x00 & ID & R & 0xA5 \\
0x01 & VERSION & R & 0x03 \\
0x02 & STATUS & R & Tabla~\ref{tab:app_status} \\
0x03 & CONTROL & W & órdenes, Tabla~\ref{tab:app_control}; se autolimpian \\
0x04 & CFG\_KD & R/W & exponente del divisor, $K_D = 2^{\text{valor}}$, rango
  efectivo 0--24 (los valores mayores equivalen a 24). Tras el reset vale
  24. Solo con el motor libre: una escritura con una orden en curso se
  ignora \\
0x05 & CFG\_PAG & R/W & página del búfer de captura, 0--3 \\
0x06--0x07 & RCT\_MAX & R & racha más larga vista por el \ac{RCT}, byte bajo
  primero \\
0x08--0x09 & APT\_MAX & R & cuenta más alta vista por el \ac{APT}, byte bajo
  primero \\
0x0A--0x0D & PETICIONES & R & contador de resembrado del \ac{DRBG}, 32 bits,
  byte bajo primero. Se satura en $2^{32}-1$ \\
0x0E & DIAG & R & Tabla~\ref{tab:app_diag} \\
0x10--0x1F & CLAVE & R & clave actual, 16 bytes. \textbf{Solo en modo test} \\
0x20--0x2F & RETO & W & nonce de 16 bytes para el reto-respuesta \\
0x30--0x3F & ETIQUETA & R & $\mathrm{CMAC}_{\text{CLAVE}}(\text{RETO})$, 16
  bytes; se lee como ceros mientras \texttt{etiqueta\_lista} vale 0 \\
0x40--0x5F & ALEATORIO & R & 32 bytes de la última orden \texttt{aleatorio},
  una generación del \ac{DRBG} independiente de la clave \\
0x80--0xFF & CAPTURA & R & 128 bytes de bits crudos, página CFG\_PAG.
  \textbf{Solo en modo test} \\
\bottomrule
\end{tabular}
\end{table}

\begin{table}[htbp]
\centering\small
\caption{STATUS (0x02).}
\label{tab:app_status}
\begin{tabular}{@{}cl>{\raggedright\arraybackslash}p{0.6\linewidth}@{}}
\toprule
Bit & Nombre & Significado \\
\midrule
0 & sembrado & el \ac{DRBG} tiene estado válido \\
1 & ocupado & hay una orden en curso; no mandar otra \\
2 & alarma\_rct & test de repetición disparado (fuente pegada) \\
3 & alarma\_apt & test de proporción disparado (fuente sesgada) \\
4 & clave\_valida & hay clave cargada en el autenticador \\
5 & captura\_llena & el búfer de captura tiene un trozo listo \\
6 & etiqueta\_lista & ETIQUETA corresponde al RETO actual \\
7 & modo\_test & el pin de test está activo \\
\bottomrule
\end{tabular}
\end{table}

\begin{table}[htbp]
\centering\small
\caption{DIAG (0x0E).}
\label{tab:app_diag}
\begin{tabular}{@{}cl>{\raggedright\arraybackslash}p{0.6\linewidth}@{}}
\toprule
Bit & Nombre & Significado \\
\midrule
0 & arrancado & el test de arranque ha pasado: una ventana completa del
  \ac{APT} (1024 muestras) sin alarma desde el reset o desde el último
  \texttt{borrar\_alarmas} \\
1 & rechazada & la última orden escrita en CONTROL se ignoró o, si era una
  siembra, se abortó. Se baja con la siguiente orden aceptada \\
2 & agotado & el contador de peticiones está en $2^{32}-1$;
  \texttt{generar} y \texttt{aleatorio} se rechazan hasta resembrar \\
\bottomrule
\end{tabular}
\end{table}

\begin{table}[htbp]
\centering\small
\caption{CONTROL (0x03). Una orden por escritura.}
\label{tab:app_control}
\begin{tabular}{@{}cl>{\raggedright\arraybackslash}p{0.64\linewidth}@{}}
\toprule
Bit & Orden & Qué hace \\
\midrule
0 & sembrar & recoge 384 bits crudos (256 de entropía y 128 de nonce),
  espera a que los tests de salud los hayan evaluado enteros e instancia el
  \ac{DRBG} desde cero \\
1 & generar & pide 256 bits; los 128 primeros pasan a ser la CLAVE y se
  cargan en el autenticador, el resto se descarta. No deja nada legible \\
2 & autenticar & ETIQUETA $\leftarrow \mathrm{CMAC}_{\text{CLAVE}}(\text{RETO})$ \\
3 & borrar\_alarmas & baja las alarmas de salud; reinicia la ventana de
  arranque \\
4 & capturar & llena el búfer de captura con bits crudos. Solo en modo test,
  con el motor libre y sin otra captura en marcha \\
5 & resembrar & como \texttt{sembrar}, pero sobre el estado actual; exige un
  \ac{DRBG} sembrado \\
6 & aleatorio & pide 256 bits y los deja en ALEATORIO \\
7 & borrar\_todo & zeroización: CLAVE, todo el estado del \ac{DRBG} y sus
  intermedios (también el núcleo \ac{AES}), RETO, ALEATORIO, ETIQUETA y la
  semilla recogida a cero, y la clave nula cargada en el autenticador. Se
  atiende siempre que el motor esté libre, también con alarma \\
\bottomrule
\end{tabular}
\end{table}

\paragraph{Reglas que impone el hardware.} Una orden rechazada no llega a
subir \texttt{ocupado} y levanta \texttt{DIAG.rechazada}.
\begin{itemize}
\item Con una alarma de salud levantada se rechazan \texttt{sembrar},
  \texttt{resembrar}, \texttt{generar} y \texttt{aleatorio};
  \texttt{autenticar} se sigue atendiendo, porque usa la clave ya generada
  y no consume entropía.
\item No se siembra antes de pasar el test de arranque
  (SP~800-90B~\S4.3~\cite{turan_2018_sp80090b}). La semilla solo entra en
  el \ac{DRBG} cuando se cierra sin alarma la ventana del \ac{APT} que
  contiene su último bit; una alarma o un \texttt{borrar\_alarmas} durante
  la recogida aborta la siembra y deja el estado anterior intacto. Una
  siembra tarda como mucho $384 + 1024$ bits crudos.
\item \texttt{sembrar} y \texttt{resembrar} se rechazan con
  \texttt{CFG\_KD} por debajo del mínimo (genérico \texttt{G\_KD\_MIN}, 6
  por defecto), por debajo del cual el cruce de dominio de reloj del bit
  pierde bits en silencio sin que los tests de salud puedan notarlo.
\item Una captura y una siembra se excluyen, así que el búfer de captura
  nunca contiene bits de una semilla.
\item El contador de peticiones no da la vuelta: en $2^{32}-1$ el motor se
  niega a generar hasta resembrar (\texttt{DIAG.agotado}). La política de
  resembrado por debajo de ese tope es del maestro.
\end{itemize}

\paragraph{Prioridad.} Si llegan varios bits de orden en una misma
escritura, solo se atiende el de mayor prioridad: \texttt{borrar\_todo},
después \texttt{sembrar}/\texttt{resembrar} (\texttt{resembrar} si van los
dos), \texttt{generar}, \texttt{aleatorio}, \texttt{autenticar},
\texttt{capturar}. \texttt{borrar\_alarmas} se atiende junto con
cualquiera de ellas, pero no debe compartir escritura con una siembra:
reinicia la ventana de arranque, así que la siembra se acepta y se aborta
en el ciclo siguiente.

\paragraph{Secuencia típica del maestro.}
\begin{enumerate}
\item Leer ID (0xA5) y VERSION (0x03).
\item Escribir \texttt{CFG\_KD} con el divisor de la caracterización, o
  dejar el valor de reset.
\item Sondear DIAG hasta \texttt{arrancado}; escribir \texttt{sembrar};
  sondear STATUS hasta que baje \texttt{ocupado}; comprobar
  \texttt{DIAG.rechazada} $= 0$, \texttt{sembrado} $= 1$ y sin alarmas.
\item Escribir \texttt{generar}, sondear, comprobar \texttt{DIAG.rechazada}.
  Ya hay clave.
\item Solo para provisionar en banco, en modo test: leer CLAVE.
\item Reto-respuesta: escribir RETO, escribir \texttt{autenticar}, sondear,
  leer ETIQUETA y compararla con el \ac{CMAC} que calcula el maestro.
\item Aleatorio para el maestro: escribir \texttt{aleatorio}, sondear, leer
  ALEATORIO.
\end{enumerate}

\paragraph{Modo test.} El pin \texttt{modo\_test}, sincronizado con dos
biestables, habilita la lectura de CLAVE y CAPTURA y la orden
\texttt{capturar}. En la \ac{FPGA} es un interruptor de la placa; en un
circuito integrado sería un fusible que se quema en producción.

\chapter{Vectores de prueba oficiales utilizados}
\label{app:vectors}

La Tabla~\ref{tab:app_vectors} recoge de dónde sale cada valor de
referencia que usan los bancos y cómo se cierra cada comprobación. Se usan
vectores oficiales siempre que una norma los publica. Donde la
disposición de entrada del hardware no coincide con ningún caso publicado,
los valores esperados salen de un modelo de referencia en Python que se
contrasta antes con los vectores oficiales y que está escrito de forma
independiente del \ac{VHDL}.

\begin{table}[htbp]
\centering\small
\caption{Valores de referencia y su origen. Las cuentas son las
comprobaciones superadas en la última ejecución
(\nolinkurl{results/tb_motor_top.log}, 177 en total, con 11 de 11
aserciones \acs{PSL} y 6 de 6 coberturas).}
\label{tab:app_vectors}
\begin{tabular}{@{}>{\raggedright\arraybackslash}p{0.2\linewidth}>{\raggedright\arraybackslash}p{0.4\linewidth}>{\raggedright\arraybackslash}p{0.32\linewidth}@{}}
\toprule
Bloque & Origen de los valores esperados & Dónde se comprueba \\
\midrule
\ac{AES}-128 (los dos núcleos) & FIPS~197~\cite{nist_2023_fips197},
  apéndices~B y C.1; SP~800-38A, F.1.1 ECB-AES128, tres bloques; mil
  cifrados encadenados repetidos; de nuevo el apéndice~C.1 tras
  \texttt{borrar} & \texttt{tb\_aes\_enc}, 18 comprobaciones, los dos
  núcleos con los mismos estímulos y obligados a coincidir \\
\ac{AES}-\ac{CMAC} & RFC~4493~\cite{song_2006_rfc4493} (los ejemplos de
  SP~800-38B~\cite{nist_2016_sp80038b}): los cuatro mensajes de 0, 16, 40
  y 64 bytes con la clave \texttt{2b7e1516\ldots}, y el de 16 bytes
  repetido; las subclaves $K_1$ y $K_2$ del mismo ejemplo & \texttt{tb\_aes\_cmac}, 7 comprobaciones \\
\ac{CTRDRBG}, modelo de referencia &
  \texttt{drbgtestvectors} del \ac{CAVP}~\cite{nist_cavp_drbg},
  \texttt{CTR\_DRBG\_no\_reseed.rsp} y \texttt{CTR\_DRBG\_pr\_false.rsp},
  AES-128 con y sin función de derivación: $4 \times 240 = 960$ casos &
  \nolinkurl{analysis/ctr_drbg_ref.py}, 960 de 960 \\
\ac{CTRDRBG}, \ac{RTL} & el modelo de referencia anterior, con la
  disposición de entrada del hardware (entropía de 256 bits, nonce de 128,
  sin personalización): \nolinkurl{analysis/gen_drbg_vectores.py} escribe
  \texttt{drbg\_vectores\_pkg.vhd} & \texttt{tb\_ctr\_drbg}, 18
  comprobaciones, incluidas la zeroización y la saturación del contador \\
Motor completo & la clave sale del ruido del anillo, así que el banco
  vuelca clave, reto y etiqueta y el \ac{CMAC} se recalcula con el modelo en
  Python & \nolinkurl{analysis/verifica_top.py}, 5 de 5 \\
Firmware contra el modelo del motor & \ac{CMAC} y ECDSA del mbedtls~3.6
  del objetivo, compilado para el PC; las lecturas firmadas se verifican
  con la biblioteca \texttt{cryptography} & \texttt{run\_host.py}, 0 fallos;
  \texttt{verifica\_lecturas.py} \\
Cortes de los tests de salud & fórmulas de
  SP~800-90B~\S4.4~\cite{turan_2018_sp80090b} para $H = 0{,}98$ y
  $\alpha = 2^{-20}$: \ac{RCT} $C = 22$, \ac{APT} $W = 1024$, $C = 596$ &
  \nolinkurl{analysis/cortes_salud.py}; comportamiento en
  \texttt{tb\_health\_tests}, 15 comprobaciones \\
\bottomrule
\end{tabular}
\end{table}

Los ficheros del \ac{CAVP} se guardan en el repositorio
(\nolinkurl{tests/vectors/drbg/}), así que la validación no depende de que la
web del \ac{NIST} siga igual. Las estimaciones de entropía del
Capítulo~\ref{ch:ch7_results} no son vectores de prueba: salen de la
herramienta SP~800-90B del \ac{NIST}~\cite{nist_ea_github}, ejecutada sin
cambios.

\chapter{Guía de reproducción}
\label{app:repro}

Todo lo que en este trabajo no necesita las placas se vuelve a ejecutar en
cada push mediante un flujo de integración continua
(\nolinkurl{.github/workflows/verifica.yml}) en una máquina Ubuntu limpia que
no es la del autor: siete trabajos, sin credenciales. El repositorio es
público: \url{https://github.com/mmggjj/TFM_ENCRY}. Las órdenes de abajo
son las que ejecuta el flujo, así que una ejecución en verde las
reproduce. Las salidas nunca se filtran: los scripts imprimen la salida
completa de cada herramienta y fallan ante cualquier error.

\begin{table}[htbp]
\centering\small
\caption{Lo que reproduce el flujo de integración continua.}
\label{tab:app_ci}
\begin{tabular}{@{}>{\raggedright\arraybackslash}p{0.27\linewidth}>{\raggedright\arraybackslash}p{0.38\linewidth}>{\raggedright\arraybackslash}p{0.27\linewidth}@{}}
\toprule
Trabajo & Orden & Resultado \\
\midrule
Modelo y estimadores & \texttt{python} \nolinkurl{analysis/ero_model.py},
  \texttt{jitter\_estimator.py}, \texttt{ctr\_drbg\_ref.py},
  \texttt{cortes\_salud.py}, \texttt{fuzz\_verifica\_lecturas.py},
  \texttt{verifica\_lecturas.py} sobre el log de la placa & autopruebas, 960
  casos del \ac{CAVP}, lecturas de la placa válidas \\
Bancos \ac{VHDL} & \texttt{python} \nolinkurl{rtl/run_tb.py} y después
  \texttt{python} \nolinkurl{analysis/verifica_top.py} & 7 bancos, informe \acs{PSL}
  como artefacto \\
Síntesis & \texttt{python} \nolinkurl{synth/sintetiza.py} \texttt{--verilog} & área por bloque en
  kGE (Capítulo~\ref{ch:ch6_asic}) \\
Firmware en el PC & \texttt{python} \nolinkurl{firmware/host_check/run_host.py}
  \texttt{--mbedtls <dir>}, y otra vez con \texttt{--sanear} & 0 fallos, también con
  ASan y UBSan \\
Compilación con ESP-IDF & \texttt{idf.py set-target}, \texttt{build},
  \texttt{size}, e \texttt{idf.py -B build\_sim -DMOTOR\_SIMULADO=1 build},
  para el verificador (ESP32) y el medidor (ESP32-C3) & 4 imágenes \\
Memoria & \texttt{tectonic tfm\_main.tex} y \texttt{tfm\_main\_es.tex} &
  los dos PDF como artefactos \\
\bottomrule
\end{tabular}
\end{table}

\paragraph{Herramientas.} Python~3.12 con \texttt{numpy}, \texttt{scipy} y
\texttt{cryptography}; GHDL de los paquetes de Ubuntu (5.0.1 con el motor
mcode en la máquina del autor) para los bancos, que se ejecutan en
\acs{VHDL}-2008; Yosys mediante \texttt{yowasp-yosys}, o la OSS CAD Suite
con el plugin de GHDL, y la biblioteca de celdas estándar IHP SG13G2 del
\ac{PDK} abierto~\cite{ihp_openpdk_repo}; mbedtls~3.6.2 del tarball de
publicación (un clon de git no trae los fuentes generados); ESP-IDF~5.3.5,
en la integración continua con el contenedor \nolinkurl{espressif/idf:v5.3.5};
tectonic~0.15. La versión de desarrollo GHDL~7.0 con el motor LLVM se cae
con las propiedades \acs{PSL} de \texttt{tb\_motor\_top} que usan nombres
externos, y no debe usarse.

\paragraph{Fuera del flujo.} Tres partes necesitan software o hardware que
el flujo no tiene:
\begin{itemize}
\item \emph{Anillo a nivel de transistor.} ngspice~47~\cite{ngspice_47} con
  OpenVAF-Reloaded~\cite{openvaf_reloaded} para compilar los modelos PSP
  del \ac{PDK}; \texttt{python} \nolinkurl{analysis/anillo_ihp_repeticiones.py}, con
  \texttt{--esquina}, \texttt{--temp}, \texttt{--vdd} y \texttt{--mismatch}
  para esquinas, alimentación y desapareamiento. ngspice no deja fijar la
  realización del ruido, así que cada ejecución es una nueva y la
  incertidumbre sale de repetir. Una realización de 400~ns tarda unos 30
  minutos.
\item \emph{Estimaciones SP~800-90B del ESP32.} La herramienta del
  \ac{NIST}~\cite{nist_ea_github} compilada desde el fuente;
  \texttt{python} \nolinkurl{analysis/entropia_90b.py} y después
  \texttt{resumen\_90b.py}, sobre las capturas guardadas en
  \nolinkurl{results/}; la prueba de reinicios con \texttt{reinicios\_90b.py}.
  Las capturas nuevas necesitan un ESP32 y \nolinkurl{analysis/campana_rng.py}.
\item \emph{\ac{FPGA}.} Vivado (nivel gratuito, con el Artix-7 incluido) y
  \nolinkurl{synth/basys3/build.tcl}; la instalación y la grabación se
  describen en \nolinkurl{docs/instalacion_vivado_windows.md}.
\end{itemize}

\paragraph{Primera campaña del ESP32.} El análisis de la primera campaña
(Secciones~\ref{sec:esp32_sp80022} a~\ref{sec:esp32_harvester}) se vuelve a
ejecutar sin la placa a partir de los ficheros guardados en
\nolinkurl{results/esp32_rng_agosto}, con \texttt{numpy}, \texttt{scipy},
\texttt{matplotlib} y \texttt{cryptography}:
\begin{itemize}
\item \texttt{python} \nolinkurl{analysis/esp32_rng/nist_sp800_22.py}
  ejecuta la autocomprobación de la batería SP~800-22 (datos del sistema
  operativo, ceros, patrón alternado).
\item \texttt{python} \nolinkurl{analysis/esp32_rng/harvester_from_sources.py}
  escribe las 1000 semillas de la validación fuera de línea en
  \texttt{seeds\_rep1.bin}, a partir de \texttt{puf\_boot1.bin} y
  \texttt{condD\_rep1.bin}.
\item \texttt{python} \nolinkurl{analysis/esp32_rng/run_campaign.py}
  ejecuta \texttt{nist\_analysis.py}, \texttt{entropy\_analysis.py},
  \texttt{temperature\_correlation.py}, \texttt{harvester\_validation.py}
  y el modo fuera de línea de \texttt{mqtt\_tls\_demo.py}. Cada uno escribe
  sus CSV y figuras en el directorio de datos, y \texttt{run\_campaign.py}
  copia las figuras en \nolinkurl{memoria/figures} con sus propios
  nombres; la memoria solo usa \texttt{nist\_heatmap.pdf}, como
  \texttt{esp32\_nist\_heatmap.pdf}. Los cuatro scripts también se ejecutan
  uno a uno, con \texttt{-d <directorio de datos>} y
  \texttt{--figdir <directorio>}. Ejecutados de nuevo el 2026-10-31 sobre
  una copia de los datos, reproducen los CSV guardados, valor a valor.
\item \texttt{python} \nolinkurl{analysis/esp32_rng/make_result_tables.py}
  escribe las tablas de resultados en \LaTeX{} en
  \texttt{tablas\_generadas.tex}, y \texttt{python}
  \nolinkurl{analysis/esp32_rng/figuras_memoria.py} dibuja las
  Figuras~\ref{fig:esp32_metrics}, \ref{fig:esp32_temp}
  y~\ref{fig:esp32_harvester} en \nolinkurl{memoria/figures}.
\end{itemize}
Una captura nueva necesita un ESP32 con \nolinkurl{firmware/esp32_rng}
compilado y grabado (\texttt{idf.py set-target esp32}, \texttt{build},
\texttt{-p <puerto> flash}; ESP-IDF~5.3, consola a 115\,200~baudios).
Después, \texttt{python} \nolinkurl{analysis/esp32_rng/capture_rng.py}
\texttt{--port <puerto>} \texttt{--reps <n>} reinicia la placa por RTS y
recoge $n$ secuencias válidas por condición, con los patrones de SRAM
(\texttt{puf\_boot<k>.bin}), las semillas y \texttt{capture\_log.csv}; por
defecto escribe en \nolinkurl{results/esp32_rng_agosto}, así que
\texttt{-o <directorio>} mantiene aparte los datos nuevos. Un reinicio por
RTS mantiene alimentada la SRAM y no da un patrón de encendido nuevo: para
la serie limpia, se desenchufa la placa, se vuelve a enchufar y se ejecuta
\texttt{python} \nolinkurl{analysis/esp32_rng/capture_puf.py}
\texttt{--port <puerto>} una vez por encendido, lo que añade un
\texttt{puf\_pc<N>.bin}. Las campañas posteriores usan el mismo firmware
con \nolinkurl{analysis/campana_rng.py}, que guarda el flujo serie íntegro
y numera las repeticiones a partir de \texttt{--primera-rep}.
